% GENERATED FILE. Edit canonical lessons, then run npm run generate:books.
% canonical-source-hash: fnv1a64:731ffefe162b7e4c
% canonical-lessons: LA-C244-raeda, LA-C244-platea, LA-C244-aquaeductus, LA-C244-thermae, LA-C244-amphitheatrum

\chapter{Travelling carriage, Broad street, Aqueduct, Public baths, Amphitheatre}
\label{ch:la-travelling-carriage-broad-street-aqueduct-public-baths-amphitheatre}
\hlchaptermodality{244}

\begin{chapteropening}
\textbf{By the end of this chapter:} \emph{I can say a travelling carriage, a broad street, an aqueduct, the public baths and an amphitheatre.}

Complete the last lesson of chapter 244: I can say a travelling carriage, a broad street, an aqueduct, the public baths and an amphitheatre.
\end{chapteropening}

\section[raeda, raedae]{raeda, raedae — a travelling carriage}

\label{lesson:LA-C244-raeda}



\begin{quote}
Before the new one: say the Latin for unlucky, then the Latin for blessed.
\end{quote}

\subsection*{You'll want to know: raeda, raedae}
\textbf{raeda, raedae} — \textquotedblleft{}a travelling carriage\textquotedblright{}.

\begin{cousinweb}
\textbf{raeda} was probably borrowed from Gaulish: Quintilian counts \emph{raeda} among the Gallic words in Latin (Lewis \& Short). Its Celtic root is the one for riding, a cousin of English \emph{ride}.
\end{cousinweb}

\begin{grammarlens}[title={What you've built}]
One more word for this chapter.
\end{grammarlens}

\subsection*{Your turn}
\begin{itemize}
\raggedright
  \item \emph{Say it:} \emph{raeda, raedae}
  \item \emph{Say it:} \emph{raeda, raedae}, once more
  \item \emph{Say it:} say \emph{raeda, raedae}
  \item \emph{Recall:} say the Latin for unlucky, then the Latin for blessed, then say \emph{raeda, raedae} again
\end{itemize}

\begin{tcolorbox}[breakable,colback=teal!4,colframe=teal!35!black,title={Before you move on}]
What does raeda, raedae mean? (\textquotedblleft{}A travelling carriage\textquotedblright{}.) Say it once more.
\end{tcolorbox}
\section[platēa, platēae]{platēa, platēae — a broad street, a square}

\label{lesson:LA-C244-platea}



\begin{quote}
Before the new one: say the Latin for blessed, then the Latin for a travelling carriage.
\end{quote}

\subsection*{You'll want to know: platēa, platēae}
\textbf{platēa, platēae} — \textquotedblleft{}a broad street, a square\textquotedblright{}.

\begin{cousinweb}
\textbf{platēa} was borrowed from Greek \emph{plateia}, \textquotedblleft{}a broad street\textquotedblright{} (Lewis \& Short), from \emph{platys}, \textquotedblleft{}broad\textquotedblright{}. English \emph{place} comes from it through French.
\end{cousinweb}

\begin{grammarlens}[title={What you've built}]
One more word for this chapter.
\end{grammarlens}

\subsection*{Your turn}
\begin{itemize}
\raggedright
  \item \emph{Say it:} \emph{platēa, platēae}
  \item \emph{Say it:} \emph{platēa, platēae}, once more
  \item \emph{Say it:} say \emph{platēa, platēae}
  \item \emph{Recall:} say the Latin for blessed, then the Latin for a travelling carriage, then say \emph{platēa, platēae} again
\end{itemize}

\begin{tcolorbox}[breakable,colback=teal!4,colframe=teal!35!black,title={Before you move on}]
What does platēa, platēae mean? (\textquotedblleft{}A broad street, a square\textquotedblright{}.) Say it once more.
\end{tcolorbox}
\section[aquaeductus, aquaeductūs]{aquaeductus, aquaeductūs — an aqueduct}

\label{lesson:LA-C244-aquaeductus}



\begin{quote}
Before the new one: say the Latin for a travelling carriage, then the Latin for a broad street.
\end{quote}

\subsection*{You'll want to know: aquaeductus, aquaeductūs}
\textbf{aquaeductus, aquaeductūs} — \textquotedblleft{}an aqueduct\textquotedblright{}.

\begin{cousinweb}
\textbf{aquaeductus} is \emph{aquae ductus}, \textquotedblleft{}a leading of water\textquotedblright{} (Lewis \& Short): \emph{aqua}, \textquotedblleft{}water\textquotedblright{}, and \emph{ductus}, \textquotedblleft{}a leading\textquotedblright{}, from \emph{dūcō}, \textquotedblleft{}to lead\textquotedblright{}.
\end{cousinweb}

\begin{grammarlens}[title={What you've built}]
One more word for this chapter.
\end{grammarlens}

\subsection*{Your turn}
\begin{itemize}
\raggedright
  \item \emph{Say it:} \emph{aquaeductus, aquaeductūs}
  \item \emph{Say it:} \emph{aquaeductus, aquaeductūs}, once more
  \item \emph{Say it:} say \emph{aquaeductus, aquaeductūs}
  \item \emph{Recall:} say the Latin for a travelling carriage, then the Latin for a broad street, then say \emph{aquaeductus, aquaeductūs} again
\end{itemize}

\begin{tcolorbox}[breakable,colback=teal!4,colframe=teal!35!black,title={Before you move on}]
What does aquaeductus, aquaeductūs mean? (\textquotedblleft{}An aqueduct\textquotedblright{}.) Say it once more.
\end{tcolorbox}
\section[thermae, thermārum]{thermae, thermārum — the public baths}

\label{lesson:LA-C244-thermae}



\begin{quote}
Before the new one: say the Latin for a broad street, then the Latin for an aqueduct.
\end{quote}

\subsection*{You'll want to know: thermae, thermārum}
\textbf{thermae, thermārum} — \textquotedblleft{}the public baths\textquotedblright{}.

\begin{cousinweb}
\textbf{thermae} comes from Greek \emph{thermos}, \textquotedblleft{}warm\textquotedblright{}: Lewis \& Short gloss \emph{thermae} as warm springs and warm baths.
\end{cousinweb}

\begin{grammarlens}[title={What you've built}]
One more word for this chapter.
\end{grammarlens}

\subsection*{Your turn}
\begin{itemize}
\raggedright
  \item \emph{Say it:} \emph{thermae, thermārum}
  \item \emph{Say it:} \emph{thermae, thermārum}, once more
  \item \emph{Say it:} say \emph{thermae, thermārum}
  \item \emph{Recall:} say the Latin for a broad street, then the Latin for an aqueduct, then say \emph{thermae, thermārum} again
\end{itemize}

\begin{tcolorbox}[breakable,colback=teal!4,colframe=teal!35!black,title={Before you move on}]
What does thermae, thermārum mean? (\textquotedblleft{}The public baths\textquotedblright{}.) Say it once more.
\end{tcolorbox}
\section[amphitheātrum, amphitheātrī]{amphitheātrum, amphitheātrī — an amphitheatre}

\label{lesson:LA-C244-amphitheatrum}



\begin{quote}
Before the new one: say the Latin for an aqueduct, then the Latin for the public baths.
\end{quote}

\subsection*{You'll want to know: amphitheātrum, amphitheātrī}
\textbf{amphitheātrum, amphitheātrī} — \textquotedblleft{}an amphitheatre\textquotedblright{}.

\begin{cousinweb}
\textbf{amphitheātrum} was borrowed from Greek \emph{amphitheatron}, \textquotedblleft{}a double theatre\textquotedblright{} (Lewis \& Short), from \emph{amphi}, \textquotedblleft{}around, on both sides\textquotedblright{}, and \emph{theatron}, \textquotedblleft{}a theatre\textquotedblright{} (EtymDB 2.1): seats all the way round.
\end{cousinweb}

\begin{grammarlens}[title={What you've built}]
One more word for this chapter.
\end{grammarlens}

\subsection*{Your turn}
\begin{itemize}
\raggedright
  \item \emph{Say it:} \emph{amphitheātrum, amphitheātrī}
  \item \emph{Say it:} \emph{amphitheātrum, amphitheātrī}, once more
  \item \emph{Say it:} say \emph{amphitheātrum, amphitheātrī}
  \item \emph{Recall:} say the Latin for an aqueduct, then the Latin for the public baths, then say \emph{amphitheātrum, amphitheātrī} again
\end{itemize}

\begin{tcolorbox}[breakable,colback=teal!4,colframe=teal!35!black,title={Before you move on}]
What does amphitheātrum, amphitheātrī mean? (\textquotedblleft{}An amphitheatre\textquotedblright{}.) Say it once more.
\end{tcolorbox}
